% ============================================================
%  第二讲
% ============================================================
\section{第二讲}

上次我们说明：给定 \(X\) 上一个不可约 \(\mathrm{SL}_2\)-局部系统，我们希望构造 \(\Bun_{\mathrm{PGL}_2}\) 上一个对应的层 \(\cF_\sigma\)，
并列出它应当满足的一些性质（当时尚未精确化）。这些性质包括：

\begin{itemize}
\item \(\cF_\sigma\) 应当是反常的；
\item \(\cF_\sigma|_{\Bun^{\mathrm{even}/\mathrm{odd}}_{\mathrm{PGL}_2}}\) 作为反常层应当不可约；
\item \(\cF_\sigma\) 应当是尖点的；
\item \(\cdots\)
\end{itemize}

在上述诸性质中，尖点性是「\(\sigma\) 不可约」这一条件的对应物；这是 Langlands 哲学的一部分。

\subsection{尖点性}
\srcpage{6}

\begin{definition}\label{def:ct}
考虑对应
\[
\begin{tikzcd}
& \Bun_B \arrow[dl,"p"'] \arrow[dr,"q"] & \\
\Bun_{\mathrm{PGL}_2} & & \Bun_{\Gm}
\end{tikzcd}
\]
其中 \(p\) 把正合列 \([0\to L\to\cE\to\cO\to 0]\in\Bun_B\) 送到 \(\cE\in\Bun_{\mathrm{PGL}_2}\)，
\(q\) 把它送到 \(L\in\Bun_{\Gm}\)。定义\emph{常数项函子}
\[
\CT_*:=q_*p^!\colon\Shv(\Bun_{\mathrm{PGL}_2})\to\Shv(\Bun_{\Gm}).
\]
对每个 \(d\in\Z\)，把 \(\Bun_B\) 换成 \(\Bun^d_B\)，类似地得到函子
\[
\CT^d_*\colon\Shv(\Bun_{\mathrm{PGL}_2})\to\Shv(\Bun^d_{\Gm}).
\]
\end{definition}

\begin{definition}\label{def:cuspidal}
称 \(F\in\Shv(\Bun_G)\) 是\emph{尖点的}，若 \(\CT_*(F)\simeq 0\)。
\yizhu{「尖点」这个译名来自自守表示：cuspidal automorphic representation 的中文惯译为「尖点自守表示」。
在几何情形，常数项函子 \(\CT_*\) 是模形式常数项 \(a_0\) 的层论类比；「尖点」即常数项为零。
本文把 ``cuspidal'' 统一译为「尖点」（作形容词时写「尖点层」），不译作「尖端的」。}{术语背景·I：本库既有译名；A：Frenkel, \emph{Langlands Correspondence for Loop Groups}, §3}
\end{definition}

\begin{remark}\label{rem:cusp-all-d}
\(F\) 是尖点的，当且仅当对任意 \(d\in\Z\)，都有 \(\CT^d_*(F)\simeq 0\)。
\end{remark}

\begin{remark}\label{rem:cusp-general}
对一般约化群，我们用所有真抛物子群及其 Levi 商来定义尖点性。
\end{remark}

回忆上一讲的图景 \eqref{fig:bunpgl2}（原文误引作 (1.3)）。下述结果表明，尖点性蕴含「在无穷远处」的消没。

\begin{proposition}\label{prop:cusp-vanish}
若 \(F\) 是尖点的，则它对 \(\Bun^d_B\) 的 \(!\)-限制在 \(d>2g-2\) 时为零。
\end{proposition}

\begin{proof}[证明梗概]
考虑映射 \(q_d\colon\Bun^d_B\to\Bun^d_{\Gm}\)。若 \(d>2g-2\)，则 \(\Bun^d_B\) 中任一点对应的正合列
\(0\to L\to B\to\cO\to 0\) 都分裂\yizhu{原文此处写 \(0\to L\to B\to\cO\to0\)，其中 \(B\) 是 \(\cE\) 之误；同页下文与脚注均用 \(\cE\)。
译文照录原文，正确形式为 \(\cE\)。}{原文笔误·A：同页自相矛盾}。于是 \(q_d\) 在 \(L\in\Bun^d_{\Gm}\) 处的纤维可以等同于
分类叠 \(B H^1(X,L)\)。

回忆：由于 \(\Ga\) 可缩，有等价
\[
\pi^*\colon\Vect\rightleftarrows\Shv(B\Ga):\pi_*,
\]
其中 \(\pi\colon B\Ga\to\mathrm{pt}\) 是投影。类似地，由于 \(H^1(X,L)\) 可缩，函子
\[
q_{d,*}\colon\Shv(\Bun^d_B)\to\Shv(\Bun^d_{\Gm})
\]
在 \(d>2g-2\) 时是等价。于是 \(\CT^d_*(F)\simeq 0\) 蕴含 \(p^!_d(F)\simeq 0\)，这正是所要的。
\yizhu{原文正文说 \(q_d\) 的纤维是 \(B H^1(X,L)\)，脚注则说纤维可典范地等同于向量叠 \(R\Gamma(X,L)[1]\)。
脚注才是正确的导出陈述。当 \(d>2g-2\) 时，\(H^1(X,L)=0\) 而 \(H^0(X,L)\neq 0\)，故
\(R\Gamma(X,L)[1]\simeq H^0(X,L)[1]\)，纤维是 \(B H^0(X,L)\)，不是 \(B H^1(X,L)\)。原文紧接着的
「因为 \(H^1(X,L)\) 可缩」也应改为「因为 \(H^0(X,L)\) 是向量空间」：向量群的有限维表示平凡，故
\(\Shv(B H^0)\simeq\Vect\)，\(q_{d,*}\) 仍是等价。}{原文笔误·A：脚注与正文不一致；Drinfeld--Gaitsgory, \emph{Geometric constant term functor(s)}, Selecta Math. 22 (2016), §2}
\end{proof}

我们还可以定义另一版本的常数项函子。

\begin{definition}\label{def:ct!}
定义 \(\CT^!\) 为函子
\[
q_!p^*\colon\Shv(\Bun_G)\to\Shv(\Bun_T).
\]
\end{definition}

\begin{remark}\label{rem:ct-duality}
形式地说，若 \(F\) 可构造，则 \(\CT_*(F)\simeq D\CT^!D(F)\)，其中 \(D\) 是 Verdier 对偶。
\end{remark}

\begin{theorem}[Drinfeld--Gaitsgory，「第二伴随」]\label{thm:second-adjoint}
有典范等价
\[
\CT^d_*\simeq\mathrm{inv}\circ\CT^{-d}_!,
\]
其中 \(\mathrm{inv}\) 是由同构 \(\Bun^d_{\Gm}\simeq\Bun^{-d}_{\Gm}\)、\(L\mapsto L^{-1}\) 诱导的函子。
\end{theorem}

\begin{corollary}\label{cor:dual-cusp}
若 \(F\) 可构造且尖点，则 \(DF\) 也尖点。
\end{corollary}

\begin{corollary}\label{cor:star-vanish}
若 \(F\) 是尖点的，则它对 \(\Bun^d_B\) 的 \(*\)-限制在 \(d>2g-2\) 时为零。
\end{corollary}

\begin{proof}
由 Drinfeld--Gaitsgory 定理，\(\CT^{-d}_*(F)\simeq 0\) 蕴含 \(\CT^d_!(F)\simeq 0\)。
与命题 \ref{prop:cusp-vanish} 的证明类似，后者蕴含 \(p^*_d(F)\simeq 0\)（\(d>2g-2\)）。
\end{proof}

\begin{corollary}\label{cor:extend}
存在开子叠 \(U_n\subset\Bun_G\)，它是有限个分层的并，使得任意尖点层 \(F\) 都能从 \(U_n\)
同时作 \(!\)-延拓与 \(*\)-延拓。
\yizhu{原文脚注 2 只说 \(F\) 从 \(U_n\)「干净地延拓」（cleanly extended）。译者按（上下文推断，无外部依据）：
「干净」应指 \(!\)-延拓与 \(*\)-延拓重合，因而延拓既不丢失也不增加信息；原文未给定义。}{原文省略·A：原文脚注 2}
\end{corollary}

\begin{remark}\label{rem:n}
事实上可以取 \(n=\max\{0,2g-2\}\)。
\end{remark}

\begin{exercise}\label{ex:second-adjoint-constant}
对 \(\Bun_G\) 上的常层验证定理 \ref{thm:second-adjoint}。
\end{exercise}

\begin{exercise}\label{ex:opposite-borel}
设 \(B^-\subset G\) 为标准反 Borel 子群，并用 \(B^-\) 代替 \(B\) 定义 \(\CT^-_*\)、\(\CT^-_!\)。
证明：定理 \ref{thm:second-adjoint} 等价于 \(\CT^-_!\simeq\CT_*\)。挑战：构造自然变换
\(\CT^-_!\to\CT_*\) 与 \(\CT_*\to\CT^-_!\)。
\end{exercise}

\begin{exercise}\label{ex:cusp-exist}
（挑战）对 \(g\ge 2\)，证明尖点层存在。
\end{exercise}

\subsection{系数}
\srcpage{7}

为了陈述对 \(\cF_\sigma\) 的其他预期，我们借用模形式理论（亦即 \(\mathrm{PGL}_2\) 的自守形式理论）中的一些动机。

回忆：权为 \(k\)、级为 \(1\) 的 holonomic 模形式
\yizhu{权 \(k\) 由 \(\sigma\) 在阿基米德位处的行为决定。}{原文省略·A：原文脚注 3}
是一个和
\[
f=\sum_{n\ge 0}a_n q^n,
\]
它在 \(|q|<1\) 时收敛，并且对任意 \(g=\begin{pmatrix}a&b\\c&d\end{pmatrix}\in\mathrm{SL}_2(\Z)\)，满足
\[
f(g\cdot\tau)=(c\tau+d)^k f(\tau),
\]
其中 \(q=\exp(2\pi i\tau)\)，\(\operatorname{Im}(\tau)>0\)。
\yizhu{原文把模形式写成 ``(holonomic) modular form''，译文照录。按上下文应为 ``holomorphic''（全纯），
``holonomic'' 是笔误（holonomic 是 \(D\)-模的术语，与模形式的全纯性无关）。}{原文笔误·A：全文术语一致性}

Langlands 猜想断言：对任意（奇、不可约、非分歧的）Galois 表示
\[
\sigma\colon\text{``}\mathrm{Gal}_{\Q}\text{''}\to\mathrm{SL}_2(\text{``}\C\text{''}),
\]
存在模形式 \(f_\sigma=\sum a_n q^n\)，使得

\begin{itemize}
\item （尖点性）常数项 \(a_0\) 为零；
\item \(a_p=\operatorname{tr}(\sigma(\mathrm{Fr}_p))\)，其中 \(\mathrm{Fr}_p\in\mathrm{Gal}_{\Q}\) 是素数 \(p\) 的 Frobenius 共轭类；
\item \(a_1=1\)，当 \((n,m)=1\) 时 \(a_{nm}=a_n a_m\)，并且
\begin{equation}\label{eq:2.1}
a_{p^{n+1}}=a_p a_{p^n}-p^{k-1}a_{p^{n-1}}.
\end{equation}
\end{itemize}

注意上述条件唯一确定 \(f_\sigma\)。

我们的目标是为 \(\Bun_{\mathrm{PGL}_2}\) 上的层提供 \(a_n\) 的类比，其中 \(a_n\) 是模形式的第 \(n\) 个 Fourier 系数。

\begin{itemize}
\item 我们把 \(n\)（\(\operatorname{Spec}\Z\) 上的有效除子）换成 \(X\) 上的有效除子 \(D\)；
\item 按照 Grothendieck 关于层--函数对应的哲学，我们把数 \(a_n(f)\) 换成向量空间 \(\coeff_D(F)\in\Vect\)。
\end{itemize}

先考虑 \(D=0\) 的情形，它类比于 \(a_1(f)\)。回忆
\[
a_n(f)=\int_{\R/\Z}f(\tau)\exp(-2\pi i n\tau)\,d\tau.
\]
这里我们把 \(\R\) 认同为 \(\mathrm{SL}_2(\R)\) 的标准 Borel 子群的单幂根。这促使我们考虑 \(\Bun_N\)，
其中 \(N\) 是 \(\mathrm{PGL}_2\) 的标准 Borel 子群的单幂根。按定义，\(\Bun_N\) 是扩张
\([0\to\cO\to\cE\to\cO\to 0]\) 的模叠。

由于一些微妙的原因，我们需要考虑一个变体
\[
\Bun^\Omega_N:=\Bun_B\times_{\Bun_{\Gm}}\{\Omega\},
\]
其中 \(\Omega\in\Bun_{\Gm}\) 是 \(X\) 上典范线丛对应的点。注意作为向量叠，有
\[
\Bun^\Omega_N\simeq R\Gamma(X,\Omega)[1].
\]
现在考虑对应
\[
\begin{tikzcd}
& \Bun^\Omega_N \arrow[dl,"p_N"'] \arrow[dr,"\phi"] & \\
\Bun_G & & H^1(X,\Omega)\simeq\A^1
\end{tikzcd}
\]
其中同构 \(H^1(X,\Omega)\simeq\A^1\) 来自 Serre 对偶。

\begin{definition}\label{def:coeff0}
定义
\[
\coeff_0(F):=\mathrm{C}^\bullet_c\bigl(\Bun^\Omega_N,\ p_N^!(F)\overset{!}{\otimes}\phi^!(\exp)\bigr),
\]
其中 \(\exp\in\Shv(\A^1)\) 是「指数层」（下面解释）。
\yizhu{\(\mathrm{C}^\bullet_c\) 是紧支集上同调。原文记作 \(\mathrm{C}^\bullet\)，未写下标；
按 \S2.2 的配对与 \(!\)-推前，此处应理解为紧支集版本。}{记号·I：本工程统一记号}
\end{definition}

\begin{example}[de Rham 情形]\label{ex:exp-dr}
取 \(\exp\) 为 \(D\)-模 \((\cO,\nabla=d-dt)\)。
\end{example}

\begin{example}[\(\ell\)-进情形]\label{ex:exp-et}
有短正合列
\[
0\to\F_p\to\Ga\xrightarrow{\ \pi\ }\Ga\to 0,
\]
其中 \(\pi\) 是 Artin--Schreier 映射 \(t\mapsto t^p-t\)。于是 \(\pi_*(\Ql)\) 带有 \(\F_p\)-作用。
固定一个非平凡特征 \(\psi\colon\F_p\to\Ql^\times\)，可以考虑 \(\pi_*(\Ql)\) 的 \(\psi\)-分量；
可以证明它是秩 \(1\) lisse 层。我们把这个层定义为 \(\exp\)。
\end{example}

\begin{remark}[Betti 情形]\label{rem:exp-betti}
在 Betti 情形，\(\exp\) 不存在；但有一些技巧可以绕开它。
\end{remark}
